\documentclass[digital]{fc-hw-template}

\author{Marcos López Merino}
\instructor{Carlos Daniel Velázquez Mendoza}
\subject{Notas de Variable Compleja I}

\usepackage[sepcounters=true,proofref={always, margin}]{phfthm}

\begin{document}
\maketitle

\section{La derivada e integral en \(\mathbb{C}\)}
\subsection{Derivada compleja}

\begin{corollary}
	Seam \(u,v \colon \Omega \subset \mathbb{C} \to \mathbb{R}\). Si \(u\) y
	\(v\) son de clase \(C^{1}\) en \(\Omega\) y satisfacen las ecuaciones de
	Cauchy-Riemann en cada punto de \(\Omega\), entonces \(f = u + iv\) es
	derivable para toda \(z \in\Omega\).
\end{corollary}

\begin{definition}
	Sea \(f \colon \Omega \subset \mathbb{C} \to \mathbb{C}\). Decimos que f es
	analítica en \(\Omega\) si \(f\) es derivable para toda \(z \in\Omega\).
\end{definition}

\subsection{Integral compleja}

\begin{definition}
	Sea \(f \colon \bigl[a, b\bigr] \to \mathbb{C}\) una función continua, la cual
	podemos escribir como \(f(t) = u(t) + iv(t)\), donde \(u(t), v(t)\)
	son funciones
	reales. Entonces la integral de \(f\) en \(\bigl[a,b\bigr]\) es el
	número complejo

	\begin{equation}
		\int_{a}^{b}\odif[sep-end=\medspace]{t} f(t) =
		\int_{a}^{b}\odif[sep-end=\medspace]{t}
		u(t) + i \int_{a}^{b}\odif[sep-end=\medspace]{t} v(t).
	\end{equation}
\end{definition}

Esta primera definición tiene las siguientes propiedades:

\begin{proposition}
	Son válidas las siguientes afirmaciones:

	\begin{enumerate}
		\item La integral es \(\mathbb{C}\)-lineal; es decir para cada
		      \(\lambda, \mu \in
		      \mathbb{C}\) y \(f,g\colon \bigl[a, b\bigr] \to \mathbb{C}\)
		      continuas se tiene

		      \begin{equation}
			      \int_{a}^{b}\odif[sep-end=\medspace]{t} \bigl(\lambda f(t) +
			      \mu g(t)\bigr) = \lambda
			      \int_{a}^{b}\odif[sep-end=\medspace]{t} f(t) + \mu
			      \int_{a}^{b}\odif[sep-end=\medspace]{t} g(t).
		      \end{equation}
		\item Para cada \(f \colon \bigl[a, b\bigr] \to \mathbb{C}\) continua se tiene

		      \begin{equation}
			      \Re \Biggl(\int_{a}^{b}\odif[sep-end=\medspace]{t} f(t)\Biggr) =
			      \int_{a}^{b}\odif[sep-end=\medspace]{t} \Re(f(t))
		      \end{equation}

		      y también

		      \begin{equation}
			      \Im \Biggl(\int_{a}^{b}\odif[sep-end=\medspace]{t} f(t)\Biggr) =
			      \int_{a}^{b}\odif[sep-end=\medspace]{t} \Im(f(t)).
		      \end{equation}
		\item Para cada \(f \colon \bigl[a,b\bigr] \to \mathbb{C}\) continua se tiene
		      \begin{equation}
			      \abs`\Bigg{\int_{a}^{b}\odif[sep-end=\medspace]{t} f(t)} \leq
			      \int_{a}^{b}\odif[sep-end=\medspace]{t} \abs{f(t)}.
		      \end{equation}
	\end{enumerate}
\end{proposition}
\end{document}
